\quellseite{396}
\chapter{§ 28. THE CHURCH BOOK OR HYMNAL}

K. Buchrucker, pp. 158–162. E. Sachsse, pp. 366–368. J. H. Schueren, pp. 105–125. J. Berndt, pp. 57–59. R. Kabisch, pp. 152–157; 247–249. A. Rude, pp. 117–122. J. Steinbeck, pp. 218–222. *W. Thilo, \emph{Das geistl. Lied i. d. ev. Volksschule Deutschlands}, ²1865. *E. Sperber, \emph{Gesch. d. Behandlung des Kirchenlieds} in Schumann u. Sperber, \emph{Gesch. d. Religionsunterrichts}, 1890. *K. Buchrucker, \emph{Der Gesangbuchunterricht}, ²1887. Schumacher, \emph{Lehrbeispiele z. Behandlung d. Kirchenliedes i. d. Volksschule}, 1885. Gattermann, \emph{50 evang. Kirchenlieder}, ³1900. Koeller, \emph{Das ev. Kirchenlied i. d. Volksschule}, 1904. O. Zuck, \emph{Das Kirchenlied}, ³1906. Schultz and Triebel, \emph{Die gebräuchlichsten Lieder d. ev. Kirche}, ¹⁷1907. Lehmensick, \emph{Kernlieder der Kirche in Stimmungsbildern}, 1907. Schlegel, \emph{Praeparationen fuer Kirchenlieder u. Psalmen}, 1908. *Achenbach, \emph{Behandlung d. Kirchenlieds auf historischer Grundlage. Lehrbeispiele nach psychologischer Methode}, ⁴1910. *F. Niebergall, \emph{Bibl. Geschichte, Katechismus, Gesangbuch}, 1910. Nic. Smith, \emph{Hymns Historically Famous}, 1901. Th. Brown and H. Butterworth, \emph{The Story of the Hymns and Tunes}, 1906. *W. L. Hunton, \emph{Favorite Hymns. Stories of the Origin, Authorship, and Use of Hymns We Love}, 1917.

Since training for participation in congregational worship is one of the Church’s duties toward the rising generation, and since the Church Book or Hymnal is the indispensable manual for worship, it will also be necessary to make the young acquainted and familiar with this book. Already in the sixteenth century when, as a rule, Church Books were not placed in the hands either of the children or adults, the catechism, for this very reason, contained as an appendix some hymns, prayers, and other parts of the liturgy, even the forms for Baptism and marriage.

Public worship adjusts itself to the seasons of the church year; and this idea of the church year is worked out in the Hymnal both in the liturgical section and in that containing the hymns. Our youth must therefore be introduced to the
\quellseite{397}
meaning of the church year. This must be done very carefully in our country where the Lutheran Church is surrounded by so many churches either altogether ignorant of the idea of the church year or acquainted with mere fragments of it. As the outstanding redemptive acts of God pass in review before our eyes in church year after church year, they invite us time and again to devout contemplation and to earnest application; neglecting the church year we would deprive ourselves of a most valuable element of religious instruction. Special lessons on the church year are, perhaps, not necessary; the information can well be imparted incidentally (in connection with the reading of the pericopes where that is still done, or in Biblical History if the stories have been selected with a view to the church year); in any event, where it has not been done previously, it should be given in a concise and coherent form in connection with the instruction of confirmands. Of course, care must be taken to avoid unnecessary historic details and popular but unwarranted theories concerning the selection of the pericopes. What has been said in our Explanation of Luther's Catechism (pp. 151–155) might be sufficient for the beginning, later to be supplemented in the Young People's Society.

The next subject requiring the attention of the catechist are the \emph{hymns}. They are the sublimate of the life of the Church; they reveal the experiences of sin and grace through which other Christians before us have passed; they witness to the abiding power of faith that has strengthened and comforted the Christian saints; they are the fairest flowers in the garden of evangelical piety. This is the very factor constituting their educational value. When the young perceive how these men of God have been sustained in affliction by their faith and by grace divine, how they have found shelter and comfort and cheer in God's omnipotent Word, their own respect for a life of faith will increase and they
\quellseite{398}
will gain greater confidence in the fountains of Christian life from which these hymns have streamed; and they will tremble at the awful reality of sin when they see how it has pressed sighs of anguish from the lips of the Christian poets. Another reason why hymns should be employed in religious instruction is that they sink more deeply into the soul when gladly and expressively sung than the spoken word; not rarely it occurs that when all else is forgotten, a hymn awakens in the soul and unfolds saving force. If there is to be fresh, vigorous congregational singing—and it is the predominance of congregational instead of choir singing that corresponds best to the Reformation principle—the young must be drilled with diligence in the melodies of those hymns; and they should be taught to sing them in their \emph{rhythmical setting} (which is finding favor owing to the efforts of Tucher, Layritz, and especially of Zahn), not in their perverted forms dating from the time of rationalism. Compare, e.g., the musical setting of the chorals in the \emph{English Hymnal} of the Missouri Synod, in the \emph{Hymnal of the Common Service Book}, in the \emph{Norwegian Lutheran Hymnary}, and in the \emph{Wartburg Hymnal for Church, School, and Home} (ed. by O. Hardwig).

In \emph{Gottes Wort im Feld und daheim}, O. Eberhard says: “The church hymn is not, like its brother, the folk song, the property of merely a certain period of life, that of youth; it belongs to the whole congregation, to all periods of life. It is learned by heart by the children, criticized by those in a state of inner ferment, explained by the lessons of life, learned anew and comprehended by the adults, fathomed more deeply through progressive experience, transfigured by age, and tested by death. The hymn is your companion from the cradle to the bier, the expression of your every need. Its deep, central notes not only thunder down from the organ loft; they also rise from the pew. And when the time for singing is past, it trembles upon the lip as a prayer; when the word loses its power in depths of woe, the old comforters once more begin to gleam; and in the darkest night of suffering they sparkle as inextinguishable stars” (\emph{Kirch. Zeitschr.} 1915. pp. 278 ff.). And Sneath, Hodges, and
\quellseite{399}
Tweedy, express themselves as follows: “One of the best gifts which parents can give to their children is constant familiarity, intelligent appreciation, and hearty and sincere use of the world's great hymns. Like all good things it will demand time and effort; but the expenditure will be more than compensated by the gain” (\emph{The Religious Training in the School and Home}, p. 163).

F. W. Herzberger says:

\begin{quote}
“Our Lutheran Church is pre-eminently the Singing Church of Evangelical Christendom. No other Church can rival her in the rich, soulful music in which she sings her immortal hymns. Countless other songs and melodies have been composed in their day, delighted their audience for a short while, and then passed into hopeless oblivion. Our majestic Lutheran chorals, however, have survived the wrecks of time, and are still today the delight of all true lovers of sacred music, irrespective of creed or language. “The Lutheran Church,” says Dr. Schaff, the noted Reformed theologian, “draws the fine arts into the service of religion, and has produced a body of hymns and chorals, which, in richness, power, and unction, surpass the hymnology of all other Churches in the world.” The late Alexander Guilmant, a Frenchman and devout Catholic, the unrivaled master of the organ in his day, declared that the Lutheran chorals are the most heart-stirring and inspiring tunes in the whole realm of sacred music. Now what is it that gives to our Lutheran chorals or church tunes their imperishable charm? Knowing their history as we do, we must say that it is the spirit of heroic faith, singing in every note its profound adoration of the merciful and omnipotent God that makes these old Lutheran chorals so universally and solemnly impressive in their character. They are alive with pure and holy devotion. They thrill the very depth of the Christian heart because they are born from the deepest and holiest passion of their inspired singers. With few exceptions, they were composed in the heroic days of the Reformation and the Thirty Years’ War, days that called for heroic courage to believe and confess the truth as it is in Jesus; days that demanded heroic submission to the inscrutable ways of our God and Redeemer. The same spirit of sublime, God-given heroism that inspired the texts of our immortal hymns also inspired their heart-stirring tunes. Hence the tunes are an integral part of our hymns. Deprive our hymns of their historic musical setting, sing them to a newer, modern tune, and you have deprived the rose of the fragrance she alone possesses, you have robbed the nightingale
\quellseite{400}
of her most rapturous note. You may then have a sorry hybrid of a poem and some sort of tune, but nevermore the original, forceful, edifying, compact hymn! For in our Lutheran hymns the text and the tune are welded together as inseparable as body and soul in man. The reason is that one and the same spirit of holy devotion gave birth to the texts as well as the chorals, or tunes, of our Lutheran hymnology. Broadly speaking then, our Lutheran chorals are pre-eminently devotional in character.

“It is different with the hymns and tunes of the eighteenth century. That was the time of decaying orthodoxy, and it witnessed the rise of Pietism in Germany and of Methodism in England. Speaking of English tunes in particular, it is a well-known fact that the Reformed Churches of Great Britain at first possessed no chorals of their own. Some of them (e.g., the Episcopal Church) originally borrowed their sacred tunes and even many hymns from the Lutheran Church of Germany. Others (e.g., the Presbyterians) contented themselves with chanting the Psalms of the Bible. They declared all “man-made” tunes and hymns to be inventions of the devil. When, however, Methodism swept over the British Islands, it produced the two greatest hymn-writers of the English-speaking world, Isaac Watts (1674—1748) and Charles Wesley (1708—1788). They were followed by other hymn-writers, both in England and America, whose songs have been set to original tunes. But what is their character? Like the emotional spirit that fostered them, they are, with a few classical exceptions, shallow, insipid, and lacking in that deep reverence of feeling that solemn harmony of tone which characterizes our old Lutheran chorals. It is true, they call themselves Gospel-hymns, but upon closer inspection you will find that very many of them contain very little Gospel and much less of true choral music. Many of them are unevangelical in text, urging and exhorting the sinner to consecrate himself to God by his own powers. Others are so silly and meaningless that sincere Christians in these churches, among them President Woodrow Wilson, have publicly protested against their further use. In their musical setting, particularly, these sensational Gospel-hymns are but little removed from the degenerate and discordant “rag-time” tunes with which the Salvation Army fills the streets of our large cities at night. How much these decadent church-tunes of a more recent date have served to vitiate the popular taste for sacred music needs no further comment.
\quellseite{401}
“We know that these old Lutheran hymns and tunes are not popular with the English-speaking people. But how can they be? Our English populace does not know them and therefore has still to learn them. And they can be learned by English people just as readily as they are learned by German or Norwegian folks. For the last ten years we are conducting a mission school in the tenement district of St. Louis, and we invite everybody and anybody to convince himself if our old Lutheran tunes cannot be learned by children of almost every nation under the sun! Again we can point to the negroes in our Colored Mission, who are originally neither German nor Norwegian, and who sing our Lutheran hymns with a vim as though they had learned them at their mother's breast. We know it takes time and patience to teach our rising English-speaking generations these noble hymns, but the joy at hearing our English youths finally singing the grand old hymns of the Reformation and glorifying their God and Redeemer in them richly repays you for all the trouble. When we hear of a Lutheran pastor who studiously avoids giving out our historic hymns and chorals in public worship, we cannot help thinking that he is either very vain and chasing after cheap popularity, or that he is very ignorant concerning the nature of true church-music, or that he is reprehensibly indolent and shirks the labor of teaching these grand tunes to his people. We hold it to be one of missionary duties of the Lutheran Church in America to acquaint the American public not only with the saving doctrine of our Church, but also with its sacred hymnology. If we Lutherans fail to do this, if we prefer the light, emotional operatic tunes of the present day to the devout, edifying tunes of our fathers, and thus suffer our historic hymns with their chorals to be forgotten, we are depriving our own posterity of the sweetest choral music this side of heaven. Therefore: ‘Lutheran Tunes for Lutheran Congregational Singing!’”\end{quote}

In order to assure intelligent singing of hymns, the teacher should read them to his class with good expression, and briefly explain them (but his explanation must not destroy the poetic perfume!). He should set forth the main thoughts, sketch the historic situation from which it took its rise, and try to enter sympathetically into the feelings of the poet. Some of the hymns, either in part or as a whole, will spontaneously fit in with other elements of religious
\quellseite{402}
instruction, especially with Biblical History and instruction in the Catechism so that from such connection the necessary light is cast upon them. For instance, in connection with the First Article, hymns such as these would be apposite: “Now thank we all our God”; “Praise to the Lord, the Almighty, the King of creation”; “All nations that on earth do dwell”; “O would, my God, that I could praise Thee”; “I sing to Thee with voice and heart”; “My soul, now bless thy Maker”; “If thou but suffer God to guide thee”; in connection with the Second Article: “All hail the power of Jesus’ name”; “Jesus, Jesus, only Jesus”; “With all my heart I love Thee, Lord”; “My dear Jesus, I’ll not leave”; “O power of love, all else transcending”; “Jesus, Thy blood and righteousness”; “Jesus, priceless treasure”; in connection with Abraham’s call: “Thy way and all sorrows”; in connection with the history of Christ’s passion: “O holy Jesus, how hast Thou offended”; “Christ, the Life of all the living”; “O sacred Head now wounded”; “Five wells I know”, etc.; compare also the following chapter. Finally a number of hymns should be committed to memory; here careful explanation is required, or else there will be confusion worse confounded. Such hymns should be the very flowers of the Church Book.

The following hymns are qualified to serve the purpose: Abide in grace, Lord Jesus; My life is hid in Jesus; Thy way and all thy sorrows; If thou but suffer God to guide thee; Now thank we all our God; Praise to the Lord, the Almighty; Good news from heav’n the angels bring; All praise to Jesus’ hallowed name; Praise God the Lord, ye sons of men; Lamb of God most holy; O sacred Head now wounded; Five wells I know; Jesus lives, no longer now; I know that my Redeemer liveth; O Holy Spirit enter in; O enter, Lord, Thy temple; Who knows how near my life’s expended; Why, my soul, thus trembling ever; God, who madest earth and heaven; Now rest beneath night’s shadows; My dear Jesus, I’ll not leave; Jesus, still lead on; Father, Son and Holy Spirit, I’m baptized in Thy dear name; A mighty fortress is our God; Take Thou my hand,
\quellseite{403}
O Father; Jesus, Lover of my Soul; Rock of Ages, cleft for me; Just as I am, without one plea; My faith looks up to Thee; Rest of the weary; The Lord my faithful Shepherd is; The Church’s one foundation. Compare § 31.

In a good Hymnal are to be found also the most important parts of the liturgy. The first task in this connection will be the introduction of the children to the \emph{structure and meaning of the common service and the communion service}. Let the teacher make clear to the pupils that we here deal with things sacred through use for the purpose of worship for over fifteen centuries; that here the adoration of God is the object in view and that we join the throng of suppliants of every time and clime. Let him make clear everything that requires elucidation, carefully practise the musical parts, arouse enthusiasm for the participation of all members of the congregation in the divine service, and especially in the liturgical part, in contrast to Rome which condemns its laity to silence. How should an active and spontaneous participation in these things be possible, if there has never been any introduction to them? From the communion service it is easy to go back to Confession and its liturgical setting. Especially the instruction of the catechumens with its pastoral character affords occasion to cast the necessary light on these acts, to discuss the salutary character of Confession and the custom of announcement for Communion, the principles of Lutheran worship, of Lutheran practice, and of Lutheran customs and Lutheran usage in general. The development of Lutheran customs and churchly usage is, on the one hand, of so much more importance in this country, as the old traditions of the home Church, in part for good reasons, have fallen into disuse; while, on the other hand, churchly usage is not only a protecting barrier for many but also a feature of healthy life in general. Among the remaining liturgical acts the baptismal ceremony requires special attention if the young at a later day, in connection with the
\quellseite{404}
duties of sponsorship, are to take an intelligent part in it. If, in connection with instruction on Baptism or Confirmation, light is cast upon the function of sponsorship with its duties and privileges, it may be possible once more to infuse life into this old institution.

\quellseite{405}
\chapter{§ 29. SECONDARY MATERIAL FOR INSTRUCTION}

A. Eckert, pp. 49. 164–175. *J. H. Schueren, pp. 97–104. O. Baumgarten, pp. 78; 82; 85–107. J. Gottschick, pp. 159–162. J. Berndt, pp. 56–67; 96. A. Rude, pp. 84–90. A. and F. Falcke, \emph{Kirchengeschichte; Praeparationen}, 1902. K. Just, \emph{Kirchengeschichtl. Unterricht}, 1903. K. Just, \emph{Kirchengeschichtl. Lesebuch}, 1903. P. Speer, \emph{Wie in unsern Schulen die Kirchengeschichte behandelt werden sollte}, 1903. F. Heyn, Kirchengeschichte in Reukauf and Heyn, \emph{Ev. Religionsbuch x 2}, 1907. Reukauf and Heyn, \emph{Kirchengeschichtliches Lesebuch}, 1907. Thraendorf and Meltzer, \emph{Kirchengeschichtliches Lesebuch}, 1907. H. Meltzer, \emph{Skizzen z. Behandlung d. Kirchengesch.}, 1909. *F. Zange, \emph{Kirchengeschichtl. Lesebuch}, 1913. Ch. F. Kent and H. B. Hunting, \emph{Witnesses for Christ}, 1913. H. W. Gates, \emph{Heroes of the Faith}, 1913. H. K. Rowe, \emph{Landmarks in Christian History}, 1914. G. H. Trabert, \emph{Church History for the People}. M. Reu, \emph{Life of Luther for the Christian Home}, 1917. C. P. Harry, \emph{Protest and Progress in the 16th Century}, 1917. *A. T. W. Steinhaeuser, \emph{Luther Primer}, 1917. K. Heilmann, \emph{Der Missionsunterricht nach Theorie und Praxis}, 1895. M. Henning, \emph{Taten Jesu in unsern Tagen}, 1906. Hemprich, \emph{Die Mission i. d. Erziehungsschule}, 1909. *G. Warneck, \emph{Die Mission i. d. Schule}, ¹⁴1911. *Th. Schaefer, \emph{Die Innere Mission i. d. Schule}, 71912. M. Hennig, \emph{Quellenbuch d. Inneren Mission}, 1912. Fred Beard, \emph{Graded Missionary Education in the Church School}. *J. F. Ohl, \emph{The Inner Mission}. *E. Pfeiffer, \emph{Mission Studies}, ²1912. *L. B. Wolf, \emph{Missionary Heroes of the Lutheran Church}. I. Boone, \emph{The Conquering Christ}, 1914. P. A. Nordell, \emph{The Modern Church}, 1914. *E. Singmaster, \emph{The Story of Lutheran Missions}, 1917. O. Koenig, \emph{Die Mission im Katechismusunterricht}, 1913. J. Richter, \emph{Die evangelischen Missionen}. J. and P. Richter, \emph{Saat und Ernte}. O. Fries, \emph{Geschichten und Bilder aus der Mission}.

Much of what has been said in connection with the Church Book constitutes material of secondary importance. Other material of secondary importance is instruction in Church History, the labors of the Church, and the constitution of the Church.

\quellseite{406}
The conviction that \emph{Church History} supplies suitable material for the instruction of our adolescent youth has been gaining ground. Of course, there is no room in such instruction for intricacies, no need of stressing dates. What is of importance is to picture the main epochs, and the men who were instruments of God in a superlative degree; above all, to describe the origin of our own Church, and of her most important symbolical book; to exhibit the sacrifices and courage of the fathers as a type for all times. What is most important of all is vividly to impress upon the young that Christ lives and reigns; that He has in His hands the threads of history, both of the Church and the world; that He will not permit His Church ever to go down in total darkness, but that He will guide her deeper and deeper into all truth through His Word; that He leads His disciples, in conformity to His own course, through grief and suffering, not to final ruin, but to eternal glory. If rightly taught, i.e., if emphasis is laid upon the biographical element, \emph{Church History} constitutes material decidedly productive of sympathetic, ethical, and religious interest, and for that reason is a great help in the attainment of the aim of religious instruction.

In point of content, such instruction readily joins itself to that of Biblical History for the Church was founded through the preaching of the gospel to Jews and Gentiles related in Biblical History. Subjects to be pictured are the spread of the Church and the persecution of Christians. The latter especially lends itself to detailed presentation, with individual pictures as object lessons. It is the history of the persecutions of the Church that brings out the power of faith to conquer the world; and at this day, it is able to arouse longing for a like faith. With the rise of Constantine, Christ celebrates His victory. Then the inner conflicts begin in the throes of which the confessions of the Church are born. Rise of monasticism. Rise of papacy. Retribution of God through the Saracens upon the lukewarm Church of Asia, Africa, and in part, of Spain. The Church among the German tribes; Charlemagne. Night in the Church, illumined by a few stars. The Holy Spirit seems to have suspended His activity; Jesus, the
\quellseite{407}
Saviour, has receded into the background, Mary and the Saints having replaced Him; instead of Christ, the pope appears to hold the sceptre; and steadily the features he exhibits become more like those of the Anti-Christ. But Christ lives. In men like Wichif and Hus, He makes the dawn to rise, with sunrise in store through the Reformation. The Spirit unfolds His power; the old promises of a Church to endure through the ages are still valid. Luther’s life. Augsburg Confession. The Catechism; divine services in the vernacular; the hymn. Lutheran Church, Reformed Church, Catholic Church; main points of difference. But hard conflicts are to follow: Thirty-years’ War. Paul Gerhardt. A. H. Francke. Era of Rationalism. But the Holy Spirit is mightier than the spirit of man, and the Church is renewed in faith. Not everything indeed is true faith which claims to be such; sectarianism and fanaticism are in evidence. Now fidelity toward the confession of the fathers is required. Fidelity needed especially in our land in which, through Muhlenberg, the Lutheran Church was founded. The Church which does not aim at outward splendor but clings to the Word in faith, is blessed by Christ; hers is the final victory.

These hints may serve as a meager outline for a course in Church History so far as the young have the right to look to the Church for information on the subject. The whole is to be pervaded by the principle: instruction for the purpose of training. For this reason it should indeed always be interesting, but never should it be sensational. As required by the principle of association, instruction in Church History should be joined to religious instruction in general, to the hymn especially. If the heroic faith of the martyrs has been the topic for discussion, the hymn is appropriate: “Onward, Christian soldiers!” If we have seen the fathers confess their faith at Nicea, the pupils may be directed to sing: “We all believe in one true God.” If the class has dwelt upon the tribulations of the Waldenses, the hymn will prove appropriate: “Thou little flock, be not afraid.” When the pupils surround in spirit the stake at which Hus is burning, let the hymn be sung: “Awake, thou Spirit who didst fire.” When the story of the Reformation flashes the Castle Church
\quellseite{408}
of Wittenberg into view, let the hymn be sounded: “Dear Christians, one and all rejoice.” The diets of Worms and Augsburg will call for their share of interest. Let it be signalled by calling for the battle hymn of the Reformation: “A mighty fortress is our God.” The story of Paul Gerhardt suggests the hymn: “Thy way and all thy sorrows.” Other hymns will suggest themselves as the occasion arises. Our \emph{Explanation of Luther’s Catechism} contains a few of the more salient features of Church History (pp. 156–167) which are merely intended to serve as basis for a mere detailed presentation to be given later. Cf. J. H. Schueren, \emph{Gedanken}, etc., pp. 107 f.

Of the various \emph{tasks of the Church}, Foreign and Inner Missions must be stressed especially. How much can be accomplished for the cause of Missions by arousing missionary enthusiasm in the young, has been shown preeminently by Herrnhut. The object must be to exhibit, through pictures from life, the desolation of the Christ-less world, and the marvelous changes wrought by the grace of God. Among the several spheres of Inner Missions, it is especially the Home Mission work, the deaconess cause, and similar works of mercy which require attention. An effective method would be to picture the lives of men who have achieved something extraordinary in those spheres; e.g., the life of Harms in the sphere of Foreign Missions, of Fliedner or Loehe in the sphere of the deaconess cause, that of A. H. Francke, or Wichern, or Geo. Mueller in the sphere of the orphan cause; or the young may be led in spirit into a deaconess home, orphans’ home, etc. Hennig’s Taten Jesu in unsern Tagen, Ohl’s Inner Mission, etc. can be of service in addition to the main factor, personal acquaintance. The labors of one’s own Church ought always to occupy the foreground; and here again those of the particular body with which one is affiliated, not merely because the brethren in the faith are first entitled to
\quellseite{409}
sympathy (Gal. 6:10), and because confederation with other Churches does not bring permanent blessing, but especially because the members of a communion should become so thoroughly identified with some specific task that they are able to say: “This is our work.” Then they will commence to work and to pray for the cause.

So far as the constitution is concerned, the chief elements of the synodical constitution can be supplied in connection with the history of one’s own synod, and the government of the local congregation may be discussed on basis of the congregational constitution.

\quellseite{410}
\chapter[§ 30. THE EDUCATIONAL AGENCIES]{§ 30. THE EDUCATIONAL AGENCIES OF THE CHURCH}

K. Buchrucker, pp. 66–115. *F. Zange, pp. 14–27. O. Baumgarten, pp. 50–85. J. Berndt, pp. 29–33. R. Kabisch, pp. 6–15. E. Chr. Achelis, pp. 338–343. H. Keferstein, \emph{Religionsunterricht u. Erziehg. zur Religion}, 1892. F. Adler, \emph{The Moral Instruction of Children}, 1892. J. McCann, \emph{The Making of Character}, 1900. Ch. W. Rishell, \emph{The Child as God's Child}, 1904. G. N. Coe, \emph{Education in Religion and Morals}, 1904. G. E. Dawson, \emph{The Child and His Religion}, 1909. E. P. St. John, \emph{Child Nature and Child Nurture}, 1911. G. Hodges, \emph{The Training of Children in Religion}, 1911. *Theo. Schmauk, \emph{In Mother's Arms}. Ch. H. Heathcote, \emph{Elements of Religious Education}, 1914. *H. Sneath, G. Hodges, and H. T. Tweedy, \emph{Religious Training in the School and Home}, 1917.

On Home Education: K. Loewe, \emph{Wie erziehen und belehren wir unsere Kinder waehrend der Schuljahre?} 1899. K. Loewe, \emph{Wie erziehe und belehre ich mein Kind bis zum 6. Lebensjahre?} ²1904. A. Matthias, \emph{Wie erziehen wir unsern Sohn Benjamin?} ⁶1907. R. Kabisch, \emph{Das neue Geschlecht}, 1913. E. Lyttleton, \emph{The Cornerstone of Education}, 1914.

On Parochial Schools compare ch. 20. On Sunday Schools compare ch. 20 and add: R. Emlein, \emph{Der Kindergottesdienst}, 1914. J. Steinbeck, pp. 224–238. P. Zauleck, \emph{Theorie und Praxis des Kindergottesdienstes in Vortraegen}. 1914. Ch. S. Albert, The School and the Church. A. H. Smith, The Lutheran Church and Child Nurture.

On Instruction of Catechumens: R. A. Kohlrausch, Der Konfirmandenunterricht, 1898. E. Simons, Konfirmation und Konfirmandenunterricht, 1900. E. Chr. Achelis, \emph{Die Bestrebungen zur Reform der Konfirmationspraxis und des Konfirmandenunterrichts} in Theol. Rundschau, 1901. Freie kirchlich-soziale Konferenz, xv and xvi, 1901. M. Gebhardt, Moderner Religions- und Konfirmandenunterricht, 1906. W. Bornemann, Der Konfirmandenunterricht u. der Religionsunterricht in der Schule in ihrem gegenseitigen Verhaeltnis, 1907. *O. Hardeland, ⁵² Konfirmandenstunden, ⁴1910. R. Steinmetz, Die Bereitung zur Konfirmation in Lehre und Leitung, 1910. O. Lorenz, Der Konfirmandenunterricht, ²1911. *K. Vogel, Seelsorgerlicher Konfirmandenunterricht, 1911. J. Steinbeck, Der Konfirmandenunterricht nach Stoffwahl, Charakter und Aufbau, ²1913.

On Instruction of the Confirmed: *H. Beck, Die Kirchliche Katechisation, 1905. *E.
\quellseite{411}
Siedel, \emph{Lebensphilosophie fuer Juenglinge}, ⁴1896. Mielke, \emph{Die religioese Fortbildung der schulentlassenen Jugend}, 1908. *E. Siedel, \emph{Der Weg zur ewigen Jugend}, ²²1909. *E. Siedel, \emph{Der Weg zur ewigen Schoenheit}, ²¹1911. A. Luettke, \emph{Unterredungen mit der konfirmierten Jugend}, ²1912. M. Reu, \emph{Die Nachpflege der Konfirmierten} in Kirch. Zeitschr., 1912. K. Knoke, \emph{Recht und Pflicht der Kirche hinsichtlich der Unterweisung der konfirmierten Jugend}, 1912. J. Schaller, \emph{The Book of Books}, 1918. M. Reu, \emph{Wartburg Lesson Helps, Senior Department}, 1918 ff. ON YOUNG PEOPLE’S SOCIETIES: Schwanbeck, \emph{Die Juenglings- und Jugendvereine}, 1890. K. Krummacher, \emph{Die evangelischen Jugendvereine und verwandte Bestrebungen}, ²1895. *P. Wurster, \emph{Die Lehre von der Inneren Mission}, 1895, pp. 380 ff. G. Hoezel, \emph{Die kirchliche Vereinsarbeit}, 1906. U. v. Hassel, \emph{Wer traegt die Schuld?} 1907. Hasse, \emph{Leitfaden fuer weibliche Jugendpflege}, 1910 f. H. Weicker, \emph{Der Jugendverein}, 1911. J. Schoell, \emph{Evangelische Gemeindepflege}, 1911. J. Eger and L. Heitmann, \emph{Die Entwicklungsjahre}, 1912 ff. Page, \emph{Evangelische Jugendpflege}: 1914. J. Steinbeck, pp. 280–315. E. Pfennigsdorf, \emph{Christus im modernen Geistesleben}, ¹⁴1914. Luther League Handbook and L. L. Topics. F. G. Detweiler, \emph{Baptist Young People at work}. W. D. Murray, \emph{Principles of the Organization of the Young Men’s Christian Association}, 1910. C. A. Barbour, \emph{Principles and Methods of Religious Work for Men and Boys}, 1912. R. C. Morse, \emph{History of the North American Y.M.C.A.}, 1913. E. H. Engelbrecht, \emph{Manual for Y.P.S.}, 1920. Augustana Luther League Council, \emph{Manual for Luther Leagues}, 1925.

But what agencies are at the disposal of the Church for the inculcation of this material?

The first agency for this purpose is the \emph{Christian home}. Although the home can be called an educational agency only in a wider sense, its significance is of greatest moment nevertheless. Not only is the home a constant factor in the educational work performed by the other agencies, but its own distinctive educational function also must continue after those of the other agencies have set in. In the home the foundation is laid for the training of the child; the home, likewise, must be the faithful custodian of that which is imparted by the other agencies to the adolescent youth. The home, accordingly, has an indispensable place in the group of educational
\quellseite{412}
agencies for the Church’s youth. He who has at heart the training of the young will not fail to give due consideration to the home whence the Church derives the material to be molded. He will rather arouse in the home an understanding for the true interests of the adolescent youth; he will remind it of its duty—doubly sacred by reason of infant Baptism—of rearing the young in the nurture and admonition of the Lord; he will invite it to co-operation with himself; and he will come to its aid with practical suggestions, counsel, and his own example.

McCunn does not exaggerate when he says in The Making of Character: “The vital matter in moral and religious education is the home as it normally is in its habitual preferences, its predominant interests, its settled estimates of persons and pursuits, its ordinary circle of associates, its standard of living, its accepted ideals of work and of amusement. For it is not only from the family but with the family eyes, that we all begin to look out upon the world. And if this first outlook is to see things for which men live in something like their true perspective, and not as distorted through the deluding medium of the home that is idle, frivolous, sordid, grasping, quarrelsome, or sentimental, this will be due far less to what is done of express educational design, far more to the ideal of life which the family consistently embodies. For it is only thus that the scale of moral valuation which the family has wrought into its life, will be likely, as the years go round, to reflect itself in the habitual feelings, estimates, and activities of its members. This kind of influence is, moreover peculiarly effective because it is made easier by the tie of natural affection. Without this, and the trustful confidence which goes with it, comparatively little can be done. And many a parent in whom the qualities which win it have been lacking, even though he may have been masterful and reasonable, has been compelled to realize his impotence. Yet, normally, the parent has a manifest advantage. That confidence which a stranger has to gain with difficulty, he finds either ready to hand, or at most less hard to win. This is double gain. It prompts a spontaneous trustfulness which opens the way for influence, and, as lesser adjunct, it invests a father’s or a mother’s disapprobation with a power to restrain and chasten such as cannot be found when love and trust are absent. In this the family is pre-eminent. No teacher, however kindly, no public authority, however

\quellseite{413}
paternal and mild, can rival it here. And if this be lost, whether by aloofness of parents, or wreck of family life, or by decay of the family as an institution, one of the purest springs of moral influence will be frozen at its source” (pp. 84 ff.).

Since it is the Church of God into which the mature as well as the immature youth has been grafted since Holy Baptism, she merely takes care of her own when she seeks to equip the home for constant co-operation in the work of child-training. A home that refuses to perform this task, cannot expect from the Church the privilege of Baptism for its children. Where, however, the home is willing to perform its duty, but for one reason or another, incapable of performing it properly, the Church must not only see to a better equipment of the home, but also provide an appropriate substitute. In the establishment of infant schools managed by some venerable matron or a deaconess and pervaded by the Christian spirit (hence quite different from the infant schools according to Froebel’s plan), also in the establishment of infant classes in Sunday School, the Church has made at least a beginning in this direction.

An educational agency additional to the home is the \emph{divine service}, however true it is that this is far more than a mere educational factor. Although it is the exception when children understand the sermon as a whole, they do understand many a section of it. With the sight of the worshipping congregation, with the liturgy, the power of song, the lesson conveyed by some striking picture here and there, it cannot be denied that the Church’s worship is a powerful educational factor in the life of the child. While there are biographies in which no wholesome effect is attributed to the attendance upon worship, there are others in which the contrary is asserted and permanent impressions are traced to the factor here stressed.

Rightly says Hodges in the chapter treating of “Sunday and the Children” in \emph{The Training of Children in Religion}: “On Sunday
\quellseite{414}
morning the children are to be taken to church. It is true that there is the peril of having the children disturb the service, and annoy the congregation; and there is also the peril of making the children hate the service. But if, in avoiding these dangers, the children are left at home, or are encouraged to consider the Sunday school a substitute for the church, there is a possibility that they never begin to go. Perhaps the best thing to do, in the midst of these perplexities, is to keep church-going as a special privilege and reward for good behavior until the age of reasonably steady habits. Then insist on attendance at church like attendance at school as a normal part of decent living and subject to precisely the same excuses. Commonly, children who are old enough to go to school are old enough to go to church. The two introductions to a wider experience may properly be made at the same period. Then nothing should be allowed to keep the child from church which would not validly keep him from school. Thus the habit is formed, and a solid contribution is made to the child's religious education. Indeed it depends much upon the character of the church service. When this, besides the sermon, provides responses and singing in which the children may join, then such a service is varied, with downsitting and uprising, and is constructed according to some understanding of human nature, the common nature not only of youth but of maturity. Even the children are not wearied.—If a choice must be made between the church and the Sunday school, on the ground that the two together are too long and wearisome, a wise preference will probably select the church. In that case, special attention will need to be given to the instruction of the children at home. One of the unedifying sights of our Christian Sunday is the spectacle of troops of children, dismissed from the church school, passing the church door, to spend the rest of the day idly. It is a prophecy of a coming generation of non-churchgoers. The habit of churchgoing, which counts for so much in the orderly religious life, is thus omitted. The boys and girls do not begin to go to church. Then, when they get through the Sunday school, in the midst of their teens, they are likely to turn their backs upon the whole system of organized religion. To this result, such training leads. Properly, the Sunday school should be a preparation for the church. The children should there be instructed to take part intelligently in the church services. They should learn there the words and music of the hymns which the congregation sing. There is opportunity,
\quellseite{415}
indeed, among the younger children for the singing of hymns which set forth the appropriate emotions of early childhood. But the older children are to learn the classic poetry of the hymnal just as they learn in the day school the classic poetry of the great masters. Some of it will far exceed their present experience, but no more so than the lines of Shakespeare and Milton, no more so than the Psalms. They will get enough for present use, and will store up treasure for the time to come. And in the church they will be able to take part with their elders. The Sunday school, however, is not a fair substitute for the church. It gives but little time to the great religious exercise of worship, and it makes no place for the sacramental side of the spiritual life. It does not bring the child into accord with the ancient, venerable, and universal expression of corporate devotion. It does not initiate him into that society in which he should have his membership all the rest of his life. It does not put him where he ought to be, and start him right.

“We probably underestimate both the endurance and the interest of children. No doubt, our forefathers demanded somewhat too much, with their protracted prayers and sermons and their two Sunday schools, one in the morning and the other in the afternoon. We are in danger of going into extreme over-anxiety as to the children’s comfort. They seem to get along pretty well with a good measure of time on week days at their schools. What is asked on Sunday is after all not excessive. The Sunday school lasts an hour, and the morning service rarely exceeds an hour and a half and is often shorter than that. It is true that at times most normal children rebel against both these forms of confinement in-doors. The best plan is to deal with their feeling about the church as we deal with their feelings about the school. We make them go to school, whether they like it or not. Sometimes they hate the school, but they almost always come into the right mind about it, and are grateful for the discipline. They may, on unusually pleasant Sundays, hate the church. No matter. If they are easily permitted to stay at home, they will despise it; and that is worse. A little wholesome temporary rebellion is better than contempt” (pp. 266 ff.).

A further noteworthy truth is found in what the same author says in \emph{Religious Training in the School and Home}: “The very building of the church is a very helpful factor in the moral and religious development of the child. . . The dignity and beauty of the architecture, the loftiness and spaciousness of the walls, the suggestion of pulpit and communion table, of font and altar, all beget a sense of solemnity
\quellseite{416}
and awe, of wonderment and of hushed expectancy that awakens and develops the moral obedience and the intelligent worship that are to be. . . In the poorest service the child gets something. His spirit is bathed in the awe and silenced in the hush. He visions the great and the good — or at least those whom he believes to be the great and the good—bound in penitence, standing jubilant in praise, instructed in the Being infinitely above him in wisdom and in goodness and in love. All this is educative. Where the minister has trained himself for this important part of his task, and is awake to the little faces that look up into his own, the hour becomes a memorable one. Many a restless and apparently thoughtless member of the junior congregation grows up to testify to the power which was exercised over him by hymn and sermon, little as they seemed to do for him at the time” (p. 7).

To be sure, attendance at divine service alone is not enough. If religious knowledge is to be something more than a group of disjointed fragments; if the will is to be habituated to that which is good, \emph{the school must be drawn into service} as an educative agency additional to home and divine worship. We do not think in this connection of the public school of our country. For though the attempt has been made again and again to incorporate in its curriculum moral and religious training in some form, such attempts are in opposition to the constitution of our country and to the character of the public school as determined by it. Whatever religion is tolerated in the public school is one from which the very heart has been torn, nor has it ever proved tenable in actual experience.

In several States it has been attempted to make the public school an agency for moral training. While, for instance, in North Carolina the teachers are merely requested to “encourage morality,” West Virginia has gone so far as to charge “all teachers, boards of education, and other school officers with the duty of providing that moral training for the youth of this State which will contribute to securing good behavior and manners, and furnish the State with exemplary citizens.” So far as religious instruction is concerned, in distinction from moral, the view has been entertained in some States that it might be rendered possible, at least in an elementary
\quellseite{417}
degree, by introducing the reading of the Bible in the school. The present status of this question in the Union is as follows: Bible reading is required by law in Alabama, Delaware, Florida, Georgia, Idaho, Kentucky, Maine, Massachusetts, New Jersey (only the Old Testament), Pennsylvania, and Tennessee. Bible reading is specifically permitted by law in Indiana, Iowa, Kansas, North Dakota, Oklahoma, and South Dakota (the only State which allows “unsectarian comment” on the portion read). State Supreme Court decisions favorable to the use of the Bible were handed down in Michigan and Texas (also in Ga., Ia., Kan., Ky., Me., Mass.); the decisions handed down in Nebraska, Wisconsin, and Ohio while not specifically forbidding, are generally considered adverse to Bible reading. The Bible is used in (some) schools under general terms of the law or by reason of its silence on the subject in Arkansas, California, Colorado, Connecticut, Maryland, Mississippi, Missouri, New Hampshire, New Mexico, New York, North Carolina, Oregon, Rhode Island, South Carolina, Utah, Vermont, Virginia, and West Virginia. The Bible is excluded from the schools by court decisions or by the Attorney General in Illinois, Louisiana, Minnesota, Washington (and apparently Nebr. and Wis.). The Bible is excluded by a ruling of the State Board of Education or the State Superintendent of Public Instruction in Arizona, Nevada, Wyoming, and Montana. In Indiana and North Dakota high school credits are given for the study of the Bible as literature outside of school hours. New York City furnishes religious teaching without charge outside of school hours in school buildings, skilled public school teachers giving the courses in religion. Cf. Wilbur F. Crafts, \emph{The Bible in School Plans of Many Lands}; Ch. N. Lischka, \emph{Private Schools and State Laws}, 1925; also the \emph{Report of the Commission on Christian Education} to the Quadrennial Meeting of the Federal Council of the Churches of Christ in America, 1916.

What are we to think of this? In order to substantiate more fully the opinion expressed above, we shall refer purposely to the opinions of non-Lutheran pedagogues. Especially noteworthy is the stand taken by Charles H. Heathcote in The Essentials of Religious Education. A decided advocate of religious education, he makes the following characteristic statement: “In theory, religion and education may be separated from each other; but in reality such a thought is impossible. The aim and goal of education and religion are virtually the same. The basis of true education is religion, and any effort to make education independent of religion narrows its scope,
\quellseite{418}
aim, and goal” (p. 7). Notwithstanding, he continues: “According to the interpretation of our constitution, religion cannot be taught in our public schools. . . We recognize the importance of reading the Bible in the schools, but we cannot call this exercise religious instruction, nor would we want it designated by such a term. When this reading is done with discrimination and without comment on the choice and splendid portions of Scriptures which should always be read, we believe great good can be accomplished in many ways. This reading may not be instructive in the analytic sense, but the mind of the pupil is impressed with the beauty and simplicity of God’s Word.—There are many educators and religionists who advocate the study of religion in the public schools. They make a contradistinction between religion and denominationalism. They advocate that religious instruction based on broad general terms of religious concepts, free from doctrinal, creedal, and denominational interpretations, could be put into the school curriculum. We realize there is much force and consideration given to their arguments, but we cannot see the feasibility of the plan. We believe in the broad interpretation of religion for it is to be thoroughly adaptable to all classes and conditions of humanity. When we speak of religion, we are, of course, referring to the broad principles of Christianity upon which the advocates of this theory agree, as it is the only religion which can give a positive civilization to the world. We do not believe that the introduction of religious instruction, even upon the basis of the broadest interpretation of Christian teachings, would work in practice. It is very evident that such a plan would not be an acceptable one to the Hebrew, Catholic, and a majority of the Protestants and many other forms of religious life which are represented in our public schools. All these conditions must be borne in mind in advocating this theory. At the same time it is well-nigh impossible to interpret religion on the broadest basis to eliminate every iota of denominational and doctrinal view-point. The public school is not a religious nor antireligious school, but it is a secular institution, and we want to see it remain such”. And afterward, on pp. 11 ff.: “We cannot agree with Dr. Seeley (\emph{Foundation of Education}, p. 248) in advocating the study of the Bible as a religious book. The Hebrew father would have every right to object, according to his religious belief and traditions, to the principles of the Christian religion being brought to his child. . . Religions cannot consistently be taught in our public schools.” Nothing remains therefore as bearer of religious education but home and Church.

\quellseite{419}
Similar conclusions, though from different premises and with different aims in view, are reached by G. A. Coe, in his work, \emph{Education in Religion and Morals}. After animadverting upon the introduction of the Bible in the public schools and upon the notion of teaching certain fundamental religious truths upon a biblical basis, he continues: “What then should be the next move toward improving the relation of the State schools toward religion? Without hesitation it may be said that the next move should be to induce the family and the Church consciously to assume their proper share of the responsibility for the character of the rising generation. Let us remove the beam that is in our own eyes. If I let weeds go to seed in my door-yard, they spread to my neighbor’s door-yard; but if I make my door-yard beautiful with flowers, I make it easier for my neighbor to beautify his own premises. As soon as the family and the Church are sufficiently aroused to begin to do their own duty, the public school question will grow wondrously simple. Strong purpose is contagious, and it has a remarkable way of finding methods. Our trouble is that we have not reached the point of giving ourselves to this reform. We are giving, instead, advice and criticism to the public schools, and in various ways we are hoping that organizations, methods, and schemes will do what only personal consecration can accomplish. We neglect the children in our homes; we do shilly-shally work in the Sunday school, and then shift to the State school the blame for the results! It is well, to be sure, to adopt at once feasible means of improving the State school, and depend upon it—any large and thorough improvement therein will wait until, through striving to build, each over against his own house, the churches and the homes have developed a proper educational consciousness among the people” (pp. 360 ff.).

When we speak of the school as an educational agency, we think primarily of the regular \emph{church, or parish, school}. We give it this name for the reason that the individual church, or parish, establishes and supervises it. Where the parish school is managed according to the basal idea that has given it birth, instruction will be imparted not only in religion but also in all other branches that have a place in the curriculum of a Christian public school. Such a school will preserve its unique character by teaching everything that is taught, according to the mind and in the spirit of Jesus
\quellseite{420}
Christ, and in harmony with the confession of the Church which sustains it. It will be arranged in grades, as required by the development of the children; it will be managed by teachers who have been specially trained for this work. Where the teachers understand their duty, faithfully fulfil it, and are themselves imbued with the spirit of Christ; where the pastor visits the school diligently, is in closest touch with its life, supervises all school activities, sees to it that the curriculum is carried out; where he pleads for the school in his sermon, in his pastoral work, and in the assembly of the congregation as a vital factor in congregational development, and where by such loyal co-operation of pastor and flock the congregation learns to recognize the great blessings of its school and freely entrusts the children of school age to it—there a constant stream of blessing will flow from such a school into the congregation, especially where the Christian home works hand in hand with the school. Beware of silencing the proposition to establish such a school with the phrase that it is an \emph{“exotic growth”}! This it cannot be for the very reason that, unlike the state school of Germany, it is a Church school (cf. pp. 188 f.). There is another reason: It was in actual existence in our country for a long time before there was such a thing as a public school, even apart from Lutheran congregations. The fact of its gradual extinction in the eastern part of our country, while partly due to other causes, is largely explained by a diminution of that spiritual energy and the lack of that vigorous Lutheran consciousness which should have pervaded the whole church life, the sphere of education included. Where no parish school was established in the West, or having been established, was permitted to pass away, the same causes are likely to have been in operation in one way or another. There is no doubt that the establishment and maintenance of such a school becomes more difficult where
\quellseite{421}
the maintenance of the mother tongue brought across the ocean is no longer one of its aims. But should the language which our children speak make us indifferent to the spirit in which geography, history, natural science, etc., are taught? Is it immaterial now whether these branches are taught in the spirit of the Word of God, which always discloses the wonder-working hand of God, or in the spirit of materialism and of the evolution theory, which pervades the textbooks of our public schools openly or in disguise; in the spirit which knows of no consummation of all things until the present process of development shall have been terminated by the judgment of God, or in the spirit of Spencer’s philosophy and pedagogy, according to which the history of humanity is a path constantly climbing higher—up to the highest heaven, which mankind is to create for itself by the labor it performs in its own behalf? Shall the natural man and an earth-born morality have the last word in the premises? Let us try to understand the idea in all its bearings: For six days in the week our children are to be exposed to the views of the natural man, and—perhaps—perceive a little of eternity’s breath in Sunday School on the seventh! Which is bound imperceptibly to gain the upper hand in these circumstances, especially where the home and its spirit does not rear a strong protecting wall around the soul? The spirit from the abyss, the spirit of lies, which utterly confuses all Christian concepts in the soul, and in a masterly manner represents as Christian what never can be such, is at work more busily than ever. Hence it is more necessary than ever to protect the souls of the rising generation from their earliest youth by instruction emanating from the spirit of Christ. And the \emph{Christain parish school} is the first place where that should be expected to be found. It is true that today when the public school system has reached a high degree of efficiency in equipment and method,
\quellseite{422}
the establishment and maintenance of a parish school is much more difficult than formerly, and that it offers no prospect of permanence in the cities unless several teachers are in charge; but we have quite a number of opulent city congregations which are equal to the financial demands thus made upon them, the more so as healthy co-operation shall tend to supersede in our Church the wretched competition of the present.

While it is therefore the holy duty of the Church in general and of the individual congregation in particular to establish \emph{regular parish schools} and to equip them with the requisite lesson helps and teaching force, conditions are not rare in which the establishment of \emph{regular parish schools} is precluded. Wherever that is the case, it is desirable, especially in country districts, to make all possible legitimate efforts to have \emph{Christian, Lutheran teachers} called to the public schools where—without in the least modifying the non-religious character of the school, but in a manner comporting with their Christian character—they may teach at public expense for as many months as is required by State law. Rome knows what she is doing when she sees to it that a disproportionately large number of teachers of her faith are given charge of the public schools. While the motive, which cannot be any other than the ultimate domination of the State by the Church, is questionable, the breadth of vision and the energy in the execution of the plan may serve as an example. In the months not taken up by public duty, such Lutheran school teachers, then, may teach parish school at congregational expense. — In cities where the establishment of a regular parish school should prove to be impossible, efforts should be made to send the children to a parish school, at least from the sixth to the ninth year. When they subsequently have to attend the public school, the spiritual influence of the first school years can be perpetuated
\quellseite{423}
by means of the Saturday and Sunday School. Where the distance from the parish school is so great that children from six to nine years of age are unable to cover it, such a partial parish school should be attended by the children during the years of preparation for confirmation.

Should even this be precluded by the shortsightedness of the leaders in the congregation or by a lack of training in the spirit of sacrifice and by unwillingness to be trained, efforts should be made to have the local or, still better, the State authorities \emph{set aside for instruction in religion a certain number of hours of the time claimed by the public school}. In proportion as the necessity for religious training shall impress itself upon the public conscience, such a plan will meet with the sympathy of the school authorities, unless the spirit of syncretism, now in the ascendant throughout the country, and aiming at a universal religion, shall foil the plan. Such instruction in religion by the Church should be supervised by a pastor conversant with teaching methods, and be imparted by teachers of a decidedly Christian character whose general equipment for teaching has been supplemented by a special course in religion and Christian pedagogy, arranged by the Church. The effect upon the child will be much deeper and more enduring if the teacher of religion is at least the equal of his colleagues in the public schools in point of technical training; if the teaching facilities are no less adequate than those obtaining in public instruction; if religious instruction is imparted in the rooms of the public school building, and if no hours are appointed for such instruction when the children are worn out by fatigue. Although the Christian character of the teacher is bound to be a decisive factor in religious instruction, no matter how adverse the conditions, there are psychological and pedagogic reasons that favor the measures above recommended. By
\quellseite{424}
systematic effort, such arrangements will generally, though not always, prove feasible. Cf. § 20.

Should the public authorities refuse to defer to the request of the congregation in the matter, it would not be advisable to wait for a change of sentiment. In that case, or if a parish school, either in the regular or partial form, should prove impracticable, \emph{a combination of Summer, Saturday, and Sunday School} will largely solve the problem. At least in the place where the pastor resides, this combination is feasible and has long since passed the experimental stage. It was subjected by me to a practical test as early as thirty-eight years ago; and since that time it has been found practicable in bilingual as well as exclusively English-speaking congregations. \emph{The Summer School} should be conducted for eight weeks at least three hours each day. Not a few pedagogues have stated as their opinion that a vacation of three months is accompanied by great disadvantages. Not a few local school authorities have established public summer schools; in some cities vacations have even been abolished, the time of daily instruction being shortened instead. Such measures constitute an irrefutable answer to the assertion that a religious summer school is a measure of cruelty against the children and an infringement upon the leisure to which they are entitled. \emph{Saturday School} is conducted throughout the year, with the possible exception of the eight weeks occupied by the Summer School; from three to five hours should be devoted to it. All the children of the congregation from the sixth or the eighth year respectively should attend these schools. Since only religious instruction, in its various ramifications, is given (Biblical History, Biblical Geography, Church History), much can be attained for the cause of education and training of children, provided the schools are characterized by system and conscience. Results such as may be attained in a properly conducted regular church school, however, are indeed beyond the scope of partial
\quellseite{425}
parish schools or these combination schools, because the influence of the latter is in part paralyzed by that of the non-religious State school, and because the factor of habit, which looms large in the sphere of child training, is not sufficiently operative. The Saturday School should supplement the Sunday School. The “presentation” of the Bible story given in the Sunday School is to be followed by the “penetration” and “application” in the Saturday School (cf. § 31). At any rate, the Saturday and the Sunday Schools must adjust themselves to each other. Any school whatever which aims at respectable achievements, the Summer School and the Saturday School included, should have a regular curriculum. So also \emph{the Sunday School should have a regular curriculum} embracing all grades. But such a curriculum cannot be established unless one has clear views in regard to the aim and tasks of the Sunday School. Where a regular parish school is in existence, and really attended by the whole congregational youth from the sixth year upward, the Sunday School is hardly a necessity. In that case it would be better to change it into a junior congregation, provided this would not affect attendance at the regular services, especially so far as the upper classes are concerned (cf. p. 391). The situation is different where the Sunday School can be made a missionary agency for the purpose of gathering religiously neglected children. Where there is no parish school, or where the parish school is not attended by all the children of the congregation, the Sunday School is a necessity, and its session is a school rather than a divine service. There must be teaching there, reciting, the covering of a definite amount of ground; but all this must be under the control of the paramount idea: the instruction given is religious, and the object is training; an impression is to be made not upon the intellect alone, but upon the whole man, who, through such instruction, is to be trained for something higher. The curriculum, however,
\quellseite{426}
can be carried out only where the several assistants or Sunday School teachers, regularly meet the pastor for their own instruction first, in order to discuss the subjects to be taught, both in regard to content and their bearings upon the several grades, and thereupon make thorough preparation for their classes. This involves the other proviso that, with the exception of the infant classes, the school is occupied not with diverse but with the same material (one and the same Bible story).

To the educational agencies mentioned—the home, the divine service, the church school in one form or another—must be added the \emph{instruction of catechumens}, which immediately precedes confirmation and dare never be omitted. Where there is a properly maintained regular parish school, three hours per week for six months will suffice for this instruction, especially if the pastor himself imparts religious instruction in the upper grade of the parish school. Where there is a partial parish school, so constituted that it is attended by the children of twelve and thirteen years of age for two years, the same number of hours will prove adequate. If there is neither a full nor a partial parish school nor “week day religious instruction,” catechumenal instruction should be given four times a week for two years unless incalculable injury is to be inflicted upon the congregation. If the catechumenal lesson could be given from eight to nine, instead of from four to five, a stronger effect would be achieved inasmuch as the children would be fresh and the whole day would come under the influence of the initial lesson. If the catechetical material is offered at a time of fatigue, the danger arises that it will be identified with ennui and fatigue. If such instruction is imparted as a mere appendage to public school instruction, the soul is in danger of falling a victim to the fallacy that religion is of less importance for the tasks of life than arithmetic and rhetoric. Should anyone
\quellseite{427}
even think that he can prepare his pupils for confirmation and the Christian life with a few “lectures”, he is guilty of most conscienceless conduct both from the pedagogic and the religious standpoint. The instruction of catechumens is the most important instruction; he who is not faithful here is not fit to hold the office which preaches reconciliation with God; and the more any congregation should insist upon shortening the time devoted to instruction, the more decided the pastor should be in opposition. — Should there be a second congregation in the parish, Sunday School should be maintained there regularly and, if possible, Saturday School; and wherever possible, the children of such congregation should join those of the main congregation for instruction during a term of from four to five months at least. Only self-denial and conscientious purpose can produce order and efficiency.

With confirmation, however, religious instruction is not brought to a conclusion. While an important goal has been reached, it is, after all, but preliminary to one beyond. Confirmation should be followed by further care. \emph{The Christian home} serves the purpose of helping the young to retain and exploit the treasures acquired. It has been shown above how much depends in just these years upon the example of father and mother: their mature faith, their firmness of character ought to inspire, influence, strengthen the youths’ religious life (cf. pp. 389–391). Family devotion both morning and evening; regular attendance at divine services; common reading of good literature; common pleasure and recreation; constant self-criticism and self-control in expressing judgments, in conversation, in actions; sympathy for questions engendered by youthful doubt; toleration for everything lawful; truthfulness, firmness, and freedom in noble union — these are requisites of the Christian home. It is to be regretted that the home, but too often, betrays too little
\quellseite{428}
understanding, willingness, and readiness for such ideals, and requires a vast amount of patient, toilsome training. Ofttimes it is even an enemy of such Christian training, not only by reason of its indifferent religious and moral attitude, but also because the axiom, “Confirmation is the close of religious education,” not without the fault of the Church, has been ingrained into its very life. Therefore, if the confirmed youth is to receive adequate care and the work begun is to be continued, another educational agency must take its place beside the home. We refer to the Bible Class of the Sunday School and the Young People’s Society. Attendance at Bible Class should be the self-evident duty of all the confirmed down to the dangerous sixteenth or seventeenth year (cf. p. 264 ff.). Those who have attended it during those years, are likely to do so also afterward. The Young People’s Society exists not only as an external tie for the confirmed youth, but also, and preeminently, as a means of their further mental and spiritual culture. It adopts its own constitution; the general direction is committed to the pastor (cf. \emph{Rules for Young People’s Societies}, Waverly, 1912). Either it is divided into two sections (from 13 to 16, and 17 to 21 respectively), or it is recruited from the ranks of those who have gone through the Bible Class of the Sunday School. If attendance upon the Bible Class (also upon the meetings of the lower section of the Y.P.S.) is a matter of obligation and congregational enactment, membership in the upper section is left free. Properly conducted, the Young People’s Society is an important educational agency, a source of blessing for all participants, and a society from which later the ranks of the most useful and mature members of the congregation recruit themselves. But prospect of permanence exists only where, from the very beginning, the work conforms to a preconceived plan; otherwise it will not survive the charm of novelty. It stands to reason that it should be, in part,
\quellseite{429}
a school of training for the future assumption of the burdens and duties of active membership in the congregation. The main stress, however, must be laid upon safeguarding and experiencing the things learned during the years of school, though under new aspects (cf. § 31), and always with the practical life as background. Thus viewed and conducted, the Young People’s Society requires indeed labor; but what is a faithful pastor not ready to do when he perceives that his labors can be a blessing to the congregation! Also here Luther’s word holds true: “Christ Himself shall be our reward!”

The grading of these educational agencies is precluded from the outset so far as the home, the instruction of catechumens, and the Young People’s Society (upper section) are concerned. It is indispensable in the parish school, as well as in the Saturday, Summer, and Sunday Schools. According to what principles the grading should be done, is shown in § 22 with sufficient clearness. If no more than two grades are possible, the division line will be the ninth or tenth year; if three grades are convenient, they will extend from the sixth to the eighth, from the ninth to the eleventh, and from the twelfth to the thirteenth or fourteenth year.
